\documentclass[12pt]{article}

\usepackage[utf8]{inputenc}
\usepackage[T1]{fontenc}
\usepackage[polish]{babel}
\usepackage{listings}
\usepackage{xcolor}

\lstset{
    language=C,                     % set language
    basicstyle=\ttfamily\small,     % font style
    keywordstyle=\color{blue},      % keywords
    stringstyle=\color{red},        % strings
    commentstyle=\color{green},     % comments
    numbers=left,                   % line numbers
    numberstyle=\tiny,              
    frame=single,                   % frame around code
    breaklines=true                 % automatic line breaking
}




\title{Dokumentacja implementacyjna}
\author{Mateusz Kamieniecki \\ Jan Rękas}
\date{\today}


\begin{document}

\maketitle

\tableofcontents

\section{main.c}
W funkcji main mamy wczytywanie argumentów do odpalania programu przy użyciu
getopt(). Argumenty mają domyślnie wartości:
\begin{itemize}
	\item -i: input.txt
	\item -o: . (domyslny katalog gdzie znajduje się program)
	\item -n: output
	\item -a: 1
\end{itemize}
Jak żaden z argumentów \textit{-h} oraz \textit{-b} nie jest podany to program
automatycznie uruchamia się z flagą \textit{-b} włączoną i informuje użytkownika. W
funkcji main są też odpalane funkcje do wczytywania danych z pliku, wypisywania
danych do pliku/plików, funkcje tworzące graf z plików oraz funkcje do
przeprowadzenia algorytmów.

\section{fileio.c oraz fileio.h}

Pliki \texttt{fileio.c} i \texttt{fileio.h} odpowiadają za wczytywanie grafu z pliku tekstowego
oraz zapisywanie wyników obliczeń do pliku wyjściowego w formacie tekstowym lub binarnym.

\subsection{fileio.h}

Nagłówek deklaruje cztery funkcje:

\begin{lstlisting}
Graph* load_graph_from_file(const char *filepath);
int    find_max_vertex_id(FILE *file);
void   save_layout_human (const char *filepath, GraphLayout *layout);
void   save_layout_binary(const char *filepath, GraphLayout *layout);
\end{lstlisting}

\subsection{fileio.c}

\subsubsection{find\_max\_vertex\_id}
Pomocnicza funkcja przechodzi przez cały plik metodą \texttt{fscanf}, wczytując kolejne
linie w formacie \texttt{<nazwa> <from> <to> <waga>} i wyznacza najwyższy identyfikator
wierzchołka spośród wszystkich wartości \texttt{from} i \texttt{to}.
Po zakończeniu pętli wywołuje \texttt{rewind(file)}, aby wskaźnik pliku powrócił na
jego początek --- jest to niezbędne, ponieważ plik jest następnie wczytywany po raz
drugi przez \texttt{load\_graph\_from\_file}. Funkcja zwraca \texttt{0} gdy plik jest
pusty lub równy \texttt{NULL}.

\subsubsection{load\_graph\_from\_file}
Funkcja otwiera plik o podanej ścieżce. W przypadku błędu wypisuje komunikat na
\texttt{stderr} i zwraca \texttt{NULL}.\\
Następnie wywołuje \texttt{find\_max\_vertex\_id}, aby ustalić wymaganą wielkość
struktury grafu. Jeżeli wynik wynosi \texttt{0} (pusty plik), zwalnia zasoby i zwraca
\texttt{NULL}.\\
Graf tworzony jest z \texttt{vertices\_amount + 1} wierzchołkami --- indeks \texttt{0}
pozostaje nieużywany, dzięki czemu identyfikatory wierzchołków z pliku wejściowego
(zaczynające się od \texttt{1}) odpowiadają wprost indeksom tablicy \texttt{adj}.\\
Pętla \texttt{fscanf} wczytuje kolejne krawędzie i dodaje je do grafu wywołaniem
\texttt{add\_edge}.

\subsubsection{save\_layout\_human}
Funkcja otwiera plik do zapisu. W przypadku błędu wypisuje komunikat na \texttt{stderr}
i kończy działanie.\\
Najpierw zapisuje liczbę wierzchołków (\texttt{layout->count}) w pierwszej linii pliku.\\
Dla każdego węzła w strukturze \texttt{GraphLayout} zapisuje wiersz w formacie:
\begin{center}
	\texttt{<id> <x> <y>}
\end{center}
Po zakończeniu wypisuje potwierdzenie na standardowe wyjście.

\subsubsection{save\_layout\_binary}
Funkcja save\_layout\_binary zapisuje wyniki do pliku binarnego. Dla każdego wierzchołka
zapisuje kolejno jego identyfikator jako \texttt{int} (4 bajty), współrzędną
$x$ jako \texttt{double} (8 bajtów) oraz współrzędną $y$ jako \texttt{double}
(8 bajtów), co daje łącznie 20 bajtów na wierzchołek.

\section{graph.c oraz graph.h}
\subsection{Graph.h}
Graph.h udostępnia następujące structy:
\subsubsection{struct Edge}
\begin{lstlisting}
typedef struct Edge {
    int to;
    double weight;
    char *name;
} Edge;
\end{lstlisting}
Struct Edge reprezentuje krawędź w grafie. Zawiera ona informacje o wierzchołku
docelowym, wadze krawędzi oraz nazwie krawędzi.
\subsubsection{struct Vector}
\begin{lstlisting}
typedef struct Vector {
    Edge *data;
    int size;
    int capacity;
} Vector;
\end{lstlisting}
Struct Vector reprezentuje dynamiczną tablicę krawędzi. Zawiera ona wskaźnik na
dane, aktualny rozmiar oraz pojemność tablicy.
\subsubsection{struct Graph}
\begin{lstlisting}
typedef struct Graph {
    int verticies_num;
    Vector *adj;
} Graph;
\end{lstlisting}
Struct Graph reprezentuje graf nieskierowany. Zawiera ona liczbę wierzchołków
oraz wskaźnik na tablicę wektorów, gdzie każdy wektor reprezentuje listę
sąsiedztwa dla danego wierzchołka.
\subsection{graph.c}
Graph.c udostępnia następujące funkcje:
\subsubsection{create\_graph}
\textbf{Przyjmuje:}
\begin{itemize}
	\item int verticies\_num --- liczba wierzchołków w grafie
\end{itemize}
\textbf{Zwraca:}
\begin{itemize}
	\item struct Graph --- wskaźnik na utworzony graf lub \texttt{NULL} w przypadku błędu alokacji pamięci.
\end{itemize}
Funkcja create\_graph tworzy graf o określonej liczbie wierzchołków, alokując
pamięć dla tablicy wektorów i inicjalizując każdy wektor.
\subsubsection{add\_edge}
\textbf{Przyjmuje:}
\begin{itemize}
	\item struct Graph *graph --- wskaźnik na graf, do którego dodajemy krawędź
	\item int from --- identyfikator wierzchołka początkowego
	\item int to --- identyfikator wierzchołka docelowego
	\item double weight --- waga krawędzi
	\item const char *name --- nazwa krawędzi
\end{itemize}
\textbf{Zwraca:}
\begin{itemize}
	\item void
\end{itemize}
Funkcja add\_edge dodaje krawędź do grafu, aktualizując listę sąsiedztwa dla obu wierzchołków (ponieważ graf jest nieskierowany).
\subsubsection{free\_graph}
\textbf{Przyjmuje:}
\begin{itemize}
	\item struct Graph *graph --- wskaźnik na graf do zwolnienia
\end{itemize}
\textbf{Zwraca:}
\begin{itemize}
	\item void
\end{itemize}
Funkcja free\_graph zwalnia pamięć zajmowaną przez graf, w tym pamięć dla list
sąsiedztwa i samego grafu.
\section{node.h}
Node.h udostępnia dwa struktury:

\subsubsection{struct Node}
\begin{lstlisting}
typedef struct Node {
    int id;
    double x;
    double y;
} Node;
\end{lstlisting}
Struct \texttt{Node} reprezentuje wierzchołek w grafie, zawierając informacje o jego
współrzędnych $(x, y)$ oraz unikalnym identyfikatorze \texttt{id}.

\subsubsection{struct GraphLayout}
\begin{lstlisting}
typedef struct {
    int count;
    Node *nodes;
} GraphLayout;
\end{lstlisting}
Struct \texttt{GraphLayout} przechowuje końcowe wyniki obliczeń --- tablicę \texttt{nodes}
zawierającą \texttt{count} wierzchołków wraz z wyznaczonymi współrzędnymi. Jest to typ
przekazywany do funkcji zapisu w \texttt{fileio.c} i zwracany przez funkcje algorytmów.
\section{tutte\_embedding.c oraz tutte\_embedding.h}
\subsection{tutte\_embedding.c}
\subsubsection{find\_embedding}
\begin{itemize}
	\item struct Graph *graph --- wskaźnik na graf, dla którego obliczamy embedding
	\item struct GraphLayout *layout --- wskaźnik na strukturę, do której zapiszemy wyniki
\end{itemize}
\begin{itemize}
	\item void
\end{itemize}
Funkcja jako argument przyjmuje structy Graph oraz GraphLayout, liczy za
pomocą tutte\_embeddingu koordynaty wierzchołków i zapisuje je do struktury
GraphLayout.


\section{spectral\_layout.c oraz spectral\_layout.h}
\subsection{spectral\_layout.c}
Spectral\_layout.c udostępnia następujące funkcje:
\subsubsection{compute\_spectral\_layout}
Funkcja compute\_spectral\_layout implementuje algorytm rozmieszczenia spektralnego.
Współrzędne $x$ wierzchołków wyznaczane są na podstawie wektora Fiedlera (drugi wektor
własny macierzy Laplace'a, odpowiadający najmniejszej niezerowej wartości własnej $\lambda_2$),
a współrzędne $y$ na podstawie trzeciego wektora własnego ($\lambda_3$). Funkcja wywołuje
kolejno funkcje pomocnicze i zwraca wypełnioną strukturę GraphLayout.
\subsubsection{compute\_degrees}
Funkcja compute\_degrees oblicza stopień każdego wierzchołka grafu, iterując po jego
liście sąsiedztwa i zapisując wyniki do tablicy degrees, która jest wymagana przez
pozostałe funkcje algorytmu.
\subsubsection{laplacian\_multiply}
Funkcja laplacian\_multiply oblicza iloczyn macierzy Laplace'a i wektora bez jawnego
budowania macierzy --- korzysta bezpośrednio z listy sąsiedztwa grafu oraz tablicy
stopni wierzchołków, co pozwala uniknąć alokacji tablicy rozmiaru $n^2$.
\subsubsection{find\_eigenvector}
Funkcja find\_eigenvector wyznacza kolejny wektor własny macierzy Laplace'a metodą
potęgową z deflacją. Przyjmuje tablicę wcześniej znalezionych wektorów własnych
i ortogonalizuje względem nich wynik, aby uzyskać następny w kolejności wektor własny.
\subsubsection{orthogonalize}
Funkcja orthogonalize przeprowadza ortogonalizację Grama--Schmidta danego wektora
względem przekazanej bazy wektorów, co jest wykorzystywane w trakcie metody potęgowej
do eliminacji już znalezionych wektorów własnych.
\subsubsection{normalize\_vector}
Funkcja normalize\_vector normalizuje wektor do długości jednostkowej, dzieląc każdy
jego element przez normę euklidesową. Jest wywoływana po każdej iteracji metody
potęgowej.
\subsubsection{dot\_product}
Funkcja dot\_product oblicza iloczyn skalarny dwóch wektorów. Jest funkcją pomocniczą
wykorzystywaną przez orthogonalize oraz find\_eigenvector do sprawdzania zbieżności.

\end{document}
